% © 2026 Ing. David Jaroš — CC BY-NC-ND 4.0
%
% This work is licensed under a Creative Commons Attribution-NonCommercial-NoDerivatives
% 4.0 International License (CC BY-NC-ND 4.0).
%
% License History: Earlier drafts (up to v0.3) were released under CC BY 4.0.
% From v0.4 onward, all material is released under CC BY-NC-ND 4.0 to protect
% the integrity of the theoretical work during ongoing academic development.

\documentclass[11pt,a4paper]{article}
\usepackage{amsmath, amssymb, amsthm}
\usepackage{hyperref}
\usepackage{geometry}
\geometry{margin=2.5cm}

\newtheorem{theorem}{Theorem}
\newtheorem{proposition}{Proposition}
\newtheorem{definition}{Definition}
\newtheorem{remark}{Remark}
\newtheorem{warning}{Warning}

\title{AdS vs.\ dS: Moduli Space Geometry vs.\ Physical Spacetime\\
\large Clarification of Two Distinct Geometric Structures in UBT}
\author{Ing.\ David Jaroš}
\date{2026}

\begin{document}
\maketitle

\begin{abstract}
The Unified Biquaternion Theory (UBT) involves two geometrically distinct structures
that superficially resemble each other: (1) the \emph{physical} spacetime with
de~Sitter (dS)-like cosmological expansion, and (2) the \emph{moduli space} of the
complex time parameter $\tau$, which carries the hyperbolic geometry of the upper
half-plane $\mathbb{H}$ (negative curvature, AdS-like). This document formalises
the distinction, explains why the compact $\psi$-circle does \emph{not} generate
an AdS factor in physical spacetime, and discusses the implications for holography.
All claims are labelled with their proof status.
\end{abstract}

% ======================================================================
\section{Two Separate Geometric Structures}
\label{sec:two_structures}
% ======================================================================

\begin{warning}[Potential Confusion]
The upper half-plane $\mathbb{H}$ carrying the modular parameter $\tau$ has
constant negative (hyperbolic) curvature. This is \emph{mathematically similar}
to Anti-de~Sitter (AdS) space. However, $\mathbb{H}$ is the \emph{moduli space}
of the compact imaginary direction — it is a parameter space, not a physical
spacetime slice. The two must not be conflated.
\end{warning}

The UBT geometry separates into two layers:

\begin{center}
\begin{tabular}{|l|l|l|l|}
\hline
\textbf{Structure} & \textbf{Nature} & \textbf{Geometry} & \textbf{Curvature sign} \\
\hline
Physical spacetime $M^4$ & Real sector of biquaternion metric & de~Sitter / Lorentzian & $+\Lambda$ (positive, dS-like) \\
$\psi$-circle $S^1(R_\psi)$ & Compact imaginary direction & Flat circle & zero \\
Moduli space $\mathbb{H}$ & Parameter space of $\tau_{\rm mod}$ & Hyperbolic (Poincaré metric) & $-1$ (AdS-like) \\
\hline
\end{tabular}
\end{center}

% ======================================================================
\section{Real Sector: Physical Spacetime and de~Sitter}
\label{sec:real_sector}
% ======================================================================

\begin{definition}[Real sector of UBT]
The physical spacetime metric is:
\begin{equation}
  g_{\mu\nu}(x) \;=\; \mathrm{Re}\!\left(\mathcal{G}_{\mu\nu}[\Theta]\right),
\end{equation}
where $\mathcal{G}_{\mu\nu}$ is the biquaternionic metric tensor derived from the
Θ-field. Source: \texttt{THEORY/canonical/canonical\_axioms.tex}, Axiom A4.
\end{definition}

The real projection of the UBT metric yields a Lorentzian spacetime $(M^4, g_{\mu\nu})$
satisfying the Einstein field equations (Steps~1--5 of the GR chain, proved [L1];
see \texttt{canonical/bridges/GR\_chain\_bridge.tex}).

Cosmological observations suggest $\Lambda > 0$ (positive cosmological constant),
placing the physical Universe in a de~Sitter-like expansion phase. This is consistent
with UBT: the real-sector equations admit de~Sitter solutions exactly as in standard
GR.

\begin{proposition}[dS solutions in UBT]\label{prop:dS}
De~Sitter space with metric $\mathrm{d}s^2 = -\mathrm{d}t^2 + e^{2Ht}\mathrm{d}\vec{x}^2$
and Hubble parameter $H = \sqrt{\Lambda/3}$ is a solution of the real-sector
UBT field equations for any constant phase $\phi$ of the Θ-field.

\noindent\textbf{Status:} \textbf{[L1]} — follows directly from GR recovery
(Theorem in \texttt{canonical/geometry/gr\_as\_limit.tex}).
\end{proposition}

% ======================================================================
\section{The $\psi$-Circle and Why It Does Not Generate AdS}
\label{sec:psi_circle}
% ======================================================================

\begin{proposition}[Flat geometry of $S^1_\psi$]\label{prop:flat_circle}
The compact imaginary direction $\psi \sim \psi + 2\pi R_\psi$ carries the flat
metric of a circle $S^1(R_\psi)$. Its curvature is identically zero.

\noindent\textbf{Status:} \textbf{[L0]} — the circle $S^1(R)$ has zero intrinsic
curvature for all $R > 0$ (standard differential geometry).
\end{proposition}

\begin{remark}[Why S$^1$ does not generate AdS]
An AdS spacetime $\mathrm{AdS}_{d+1}$ requires a $(d+1)$-dimensional space with
constant negative sectional curvature. A single compact circle has:
\begin{enumerate}
  \item Dimension 1 (the minimal non-trivial example for AdS requires $\geq 2$ dimensions).
  \item Zero curvature (not negative curvature).
\end{enumerate}
Therefore, $S^1_\psi \times M^4$ is topologically a cylinder and metrically
$(\text{flat circle}) \times (\text{4D Lorentzian})$. No AdS factor is generated.

The imaginary sector of UBT contributes to phase curvature (encoded in
$\mathrm{Im}(\mathcal{G}_{\mu\nu})$), but this is the \emph{imaginary part} of
the biquaternionic metric, not a separate anti-de~Sitter spacetime.
\end{remark}

% ======================================================================
\section{Moduli Space: Hyperbolic Geometry of $\mathbb{H}$}
\label{sec:moduli_hyperbolic}
% ======================================================================

The moduli space of the torus $T^2(\tau) = \mathbb{C}/(\mathbb{Z}+\tau\mathbb{Z})$
is parametrised by $\tau \in \mathbb{H}$ (the upper half-plane). The natural
$\mathrm{SL}(2,\mathbb{Z})$-invariant metric on $\mathbb{H}$ is the Poincaré metric:
\begin{equation}
  \mathrm{d}s^2_{\mathbb{H}} \;=\; \frac{\mathrm{d}(\mathrm{Re}\,\tau)^2 + \mathrm{d}(\mathrm{Im}\,\tau)^2}{(\mathrm{Im}\,\tau)^2}.
\end{equation}
This metric has constant negative sectional curvature $K = -1$. The modular
quotient $\mathbb{H}/\mathrm{SL}(2,\mathbb{Z})$ is the moduli space of the torus.

\begin{proposition}[Hyperbolic moduli space]\label{prop:hyp_moduli}
The moduli space of the UBT torus $T^2(\tau)$ is the Poincaré upper half-plane
$(\mathbb{H}, \mathrm{d}s^2_{\mathbb{H}})$ with curvature $K = -1$.
This is mathematically analogous to $\mathrm{AdS}_2$ (2D anti-de~Sitter space),
but it is a \emph{parameter space}, not a physical spacetime.

\noindent\textbf{Status:} \textbf{[L0]} — standard result from the theory of
modular curves (see, e.g., Diamond \& Shurman, \emph{A First Course in Modular Forms},
Springer 2005).
\end{proposition}

\begin{remark}[Physical interpretation of the hyperbolic moduli space]
The hyperbolic geometry of $\mathbb{H}$ has implications for the UBT modular program:
\begin{itemize}
  \item Geodesics in $\mathbb{H}$ are semicircles and vertical lines.
  \item The imaginary axis $\{\tau = iy : y > 0\}$ is a geodesic in $\mathbb{H}$.
  \item The UBT parameter $\tau_{\rm mod} = iR_\psi/R_t$ lives on this geodesic.
  \item Modular symmetries ($\mathrm{SL}(2,\mathbb{Z})$ transformations) are
    isometries of $(\mathbb{H}, \mathrm{d}s^2_{\mathbb{H}})$.
\end{itemize}
These geometric facts constrain which values of $\tau$ are modularly equivalent
and hence \emph{physically distinct} in any modular-invariant formulation.
\end{remark}

% ======================================================================
\section{Holography Implications}
\label{sec:holography}
% ======================================================================

The AdS/CFT correspondence (Maldacena 1997) relates a theory on $\mathrm{AdS}_{d+1}$
to a conformal field theory on its $d$-dimensional boundary. Superficially, the
hyperbolic moduli space $\mathbb{H}$ might suggest an AdS/CFT-type duality for UBT.
This subsection assesses the status of such a connection.

\begin{itemize}
  \item \textbf{What is present:} The moduli space $\mathbb{H}$ has the correct
    geometry ($K = -1$, boundary at $\mathrm{Im}(\tau) = 0$) to admit a holographic
    interpretation in the modular-form sense: modular forms on $\mathbb{H}$ with
    specified boundary behaviour encode spectral information about the torus.
  
  \item \textbf{What is absent:} A dynamical gravity theory on $\mathbb{H}$ as a
    physical spacetime. The Poincaré metric here is a parameter-space metric, not
    a physical metric. There is no UBT claim of a genuine bulk gravity theory in
    moduli space.
  
  \item \textbf{Speculative connection:} The Hecke eigenvalues of the UBT theta
    functions could, in principle, be related to spectrum data of a ``boundary CFT''
    on $\partial\mathbb{H}$ (the real line). This is a mathematically coherent idea
    in the framework of $p$-adic AdS/CFT (Gubser et al.\ 2017), where primes $p$
    play the role of the AdS radius. The relationship to $p = 137$ and $p = 139$
    in the UBT Hecke program is noted as a \textbf{[MC]} (motivated conjecture)
    requiring further development.
\end{itemize}

\begin{remark}[Summary of holography status]
The moduli space of UBT complex time has the geometric structure compatible with
holographic methods, but no AdS/CFT duality has been established within UBT.
The connection is a research direction, not a result.

\noindent\textbf{Status:} \textbf{[O]} — open research direction.
\end{remark}

% ======================================================================
\section{Summary and Distinctions}
\label{sec:summary}
% ======================================================================

\begin{center}
\begin{tabular}{|p{4cm}|p{4cm}|p{4cm}|}
\hline
\textbf{Property} & \textbf{Physical spacetime $M^4$} & \textbf{Moduli space $\mathbb{H}$} \\
\hline
Nature & Physical (observable) & Parameter space (mathematical) \\
Metric & $g_{\mu\nu} = \mathrm{Re}(\mathcal{G}_{\mu\nu})$ & $\mathrm{d}s^2 = (\mathrm{d}\tau_1^2 + \mathrm{d}\tau_2^2)/\tau_2^2$ \\
Curvature sign & Positive $\Lambda$ (dS) & Negative $K=-1$ (AdS-like) \\
Dimensionality & 4D Lorentzian & 2D Euclidean \\
Symmetry & Diffeomorphism group & $\mathrm{SL}(2,\mathbb{Z})$ (discrete) \\
Observables live here? & \textbf{Yes} & No \\
GR equations hold here? & \textbf{Yes} [L1] & No (not physical spacetime) \\
AdS/CFT applicable? & No (de~Sitter, not AdS) & Potentially [MC] — open research \\
\hline
\end{tabular}
\end{center}

\bigskip

\noindent\textbf{Key statement:} The compact $\psi$-circle is a flat internal
direction; the hyperbolic structure of $\mathbb{H}$ is a property of the
\emph{moduli space} of $\tau$. Neither generates AdS in physical spacetime.
Physical spacetime in UBT is de~Sitter (positive $\Lambda$), fully consistent
with cosmological observations.

\bigskip

\noindent\textbf{Related documents:}
\begin{itemize}
  \item \texttt{research/theta\_modular\_geometry.tex} — torus and $\mathbb{H}$ formalization
  \item \texttt{AUDITS/complex\_time\_modular\_audit.md} — full modular audit
  \item \texttt{canonical/geometry/gr\_as\_limit.tex} — GR and dS solutions
  \item \texttt{canonical/bridges/GR\_chain\_bridge.tex} — GR recovery chain
  \item \texttt{research/theta\_alpha\_connection.md} — modular alpha candidates
\end{itemize}

\end{document}
